\exercisesection{}

\begin{enumerate}
\def\labelenumi{\arabic{enumi})}
\item
  设 \(\Gamma\) 是一个真 \(^{1}\) 闭凸锥（在 \(\mathbf{R}^{n}\) 中）。证明映射 \((x,y)\mapsto x+y\) （\(\Gamma\times\Gamma\) 到 \(\Gamma\) 中的）是逆紧的。由此推出，\(\Gamma\) 上任意两个测度关于映射 \((x,y)\mapsto x+y\) 都是可卷积的。
\item
  设 G 为局部紧群，\(\Gamma\) 为对 G 添加一个无穷远点 \(\omega\) 后得到的紧空间。把 G 的合成律扩充到 \(\Gamma\) ，办法是规定 \(x\omega = \omega x = \omega\) （对每个 \(x \in \Gamma\) ）。对每个测度 \(\mu\) （\(\Gamma\) 上的），一方面对应 G 上的一个有界测度 \(\mu_{1}\) ，另一方面对应一个复数 \(\mu(\{\omega\})\) 。证明：若用 \(*(\text{resp.}\widehat{*})\) 表示由 G（相应地 \(\Gamma\) ）中的乘法所定义的卷积，则 \((\mu\widehat{*}\nu)_{1} = \mu_{1}*\nu_{1}\) ，并且
\end{enumerate}

\[
(\mu \widehat {*} \nu) (\omega) = \mu (\omega) \nu_ {1} (G) + \nu (\omega) \mu_ {1} (G) + \mu (\omega) \nu (\omega)
\]

对任意两个测度 \(\mu\) 与 \(\nu\) （\(\Gamma\) 上的）成立。

